% from ega4-18-fr.tex
\oldpage[IV-4]{119}

résiduel de $A$, $f\in k[T]$ le polynôme minimal de l'image de $u$ dans $B\otimes_A k$, il existe un $F\in A[T]$ dont $f$ est l'image canonique et tel que $F(u)=0$; un tel polynôme $F$ vérifie les conditions $F'(u)\notin\mathfrak n_i$ pour $1\leq i\leq r$.
\end{lemma}

L'idéal maximal $\mathfrak m$ de $A$ est l'image réciproque de chacun des $\mathfrak n_i$ (II, 6.1.10); soit $L=B\otimes_A k$, qui est une $k$-algèbre finie. Soit $\xi$ l'image de $u$ dans $L$, et soit $n$ le rang de $L$ sur $k$; le polynôme minimal $f\in k[T]$ de $\xi$ sur $k$ est donc de degré $n$, et $L$ est isomorphe à $k[T]/f.k[T]$. Si $\mathfrak n'_i=\mathfrak n_i/\mathfrak mB$, l'hypothèse entraîne que $\Spec(L)$ est étale sur $k$ aux points $\mathfrak n'_i$ ($1\leq i\leq r$) en vertu de (17.6.1, $c)$), donc $f'(\xi)\notin\mathfrak n'_i$ en vertu de (18.4.2, (ii)). Notons maintenant que $\mathfrak mB$ est contenu dans le radical de $B$; comme $1,\xi,\ldots,\xi^{n-1}$ forment une base de $L$ sur $k$, il résulte du lemme de Nakayama que $1,u,\ldots,u^{n-1}$ engendrent le $A$-module $B$, et par suite il existe un polynôme unitaire $F\in A[T]$ de degré $n$, tel que $F(u)=0$; en outre, comme $\xi$ est racine de l'image canonique de $F$ dans $k[T]$, cette image est nécessairement égale à $f$. Mais alors l'image de $F'(u)$ dans $L$ est $f'(\xi)$, et comme $f'(\xi)\notin\mathfrak n'_i$, on a $F'(u)\notin\mathfrak n_i$ pour tout $i$.

\begin{proposition}[18.4.5]
\label{IV.18.4.5}
Soient $A$ un anneau local, $k$ son corps résiduel, $B$ une $A$-algèbre finie (resp. finie et de présentation finie). On suppose en outre, ou bien que le corps $k$ est infini, ou bien que $B$ est un anneau local. Soit $n$ le rang de $L=B\otimes_A k$ sur $k$. Pour que $B$ soit une $A$-algèbre formellement non ramifiée (resp. étale), il faut et il suffit qu'il existe un polynôme unitaire séparable $F\in A[T]$ (18.4.3) tel que $B$ soit isomorphe à un quotient de $A[T]/F.A[T]$ (resp. isomorphe à $A[T]/F.A[T]$). En outre, on peut supposer que $F$ est de degré $n$ (resp. $F$ est nécessairement de degré $n$).
\end{proposition}

Les conditions sont suffisantes en vertu de (18.4.2), sans supposer $k$ infini ni $B$ local. Pour voir que les conditions sont nécessaires, notons que si $B$ est une $A$-algèbre formellement non ramifiée, $L$ est une algèbre finie et séparable sur $k$, donc composée directe d'un nombre fini d'extensions finies et séparables $k_j$ de $k$ ($1\leq j\leq r$), $k_j$ étant donc engendrée par un élément $\xi_j$, de polynôme minimal $f_j$ ($1\leq j\leq r$) (Bourbaki, \emph{Alg.}, chap. V, \textsection~11, n\textsuperscript{o}~4, prop. 4). Montrons qu'en vertu des hypothèses faites sur $k$ ou $B$, il existe un élément $\xi$ de $L$ engendrant la $k$-algèbre $L$. C'est immédiat si $B$ est local, puisque alors $r=1$. Sinon, $k$ étant supposé infini, on peut supposer que les polynômes irréductibles $f_j\in k[T]$ sont tous \emph{distincts} (en remplaçant au besoin chaque $\xi_j$ par $\xi_j+a_j$, pour un élément convenable $a_j\in k$); si l'on pose $f=f_1f_2\cdots f_r$, il est clair que $L$ est isomorphe à $k[T]/f.k[T]$ dans les deux cas envisagés, donc est engendrée par un élément $\xi$ de polynôme minimal $f\in k[T]$ de degré $n$. Si $u\in B$ est un élément d'image $\xi$ dans $L$, le lemme de Nakayama montre que les éléments $1,u,\ldots,u^{n-1}$ engendrent le $A$-module $B$; ceci montre déjà qu'il existe un polynôme unitaire $F\in A[T]$ de degré $n$ tel que $F(u)=0$, $u$ engendrant la $A$-algèbre $B$, qui est par suite isomorphe à une algèbre quotient de $A[T]/F.A[T]$; en outre $B$ est un anneau semi-local et en chacun de ses idéaux maximaux $\mathfrak n_i$, on a $F'(u)\notin\mathfrak n_i$ d'après (18.4.4), ce qui prouve que $F'(u)$ est inversible dans $B$, donc que $F$ est un polynôme \emph{séparable}. Enfin, si $B$ est une $A$-algèbre étale, $B$ étant un $A$-module plat de présentation finie (1.4.7) est un $A$-module libre, et $1,u,\ldots,u^{n-1}$ forment une base du $A$-module $B$ (Bourbaki, \emph{Alg. comm.}, chap. II, \textsection~3, n\textsuperscript{o}~2, prop. 5), autrement dit la $A$-algèbre $B$ est

